\subsection{Validation Scenarios and Use Cases}
\label{sec:validation-scenarios}

To demonstrate the platform's capability to model diverse contagion scenarios, we developed two validation use cases representing distinct epidemiological contexts: healthcare-acquired infections in a hospital setting and childhood disease transmission in a school environment. These scenarios showcase the config-driven parameterization approach and illustrate how the same simulation engine can model fundamentally different outbreak dynamics through JSON configuration alone.

\subsubsection{Hospital Outbreak Scenario}
\label{sec:hospital-scenario}

The hospital scenario (\texttt{simulation\_hospital.json}) models healthcare-acquired infection (HAI) dynamics in a multi-ward hospital facility, representing scenarios such as MRSA, C. difficile, or COVID-19 outbreaks among patients and healthcare workers. The spatial configuration includes 24 patient beds distributed across upper and lower wards, connected by a central corridor with nurse stations, as shown in Figure~\ref{fig:simulation-screenshot}.

\begin{figure}[t]
\centering
\fbox{\parbox{0.9\columnwidth}{\centering\vspace{1.5cm}\textbf{[Figure 3: Hospital simulation in progress showing infected agents (red glow), recovered agents (blue glow), and contaminated floor tiles (gradient overlay). Real-time statistics displayed in sidebar.]}\\[0.5cm]\textit{Note: Screenshot to be captured from running application at } \texttt{http://localhost:5173}\vspace{1.5cm}}}
\caption{Hospital outbreak simulation showing healthcare-acquired infection spread. Agents with red status glows are actively infected, blue glows indicate recovered/immune agents, and floor tiles display contamination levels through color gradients (green=clean, red=severe). The sidebar shows real-time epidemic statistics including total infected, recovered, and deceased counts.}
\label{fig:simulation-screenshot}
\end{figure}

The simulation parameters are tuned to reflect the unique characteristics of healthcare settings:

\begin{itemize}
    \item \textbf{Infection Radius} (100px): Moderate transmission distance reflecting close-proximity patient care activities (bedside examinations, medication administration) and shared equipment contact.
    
    \item \textbf{Infection Probability} (0.20 per second): Lower per-second transmission rate compared to school settings, reflecting the fact that not every patient-caregiver contact results in pathogen transmission. However, cumulative exposure time in hospitals is high due to extended shifts and multiple daily contacts.
    
    \item \textbf{Floor Infection Probability} (0.20): High environmental transmission rate modeling fomite-based spread from contaminated surfaces (bed rails, medical equipment, floors) characteristic of healthcare environments where pathogens like C. difficile spores persist for months.
    
    \item \textbf{Recovery Time} (20 seconds): Extended infection duration reflecting realistic disease courses (7-14 days scaled to simulation time), representing prolonged hospitalization periods for HAI patients.
    
    \item \textbf{Fatality Probability} (0.08): Elevated mortality rate (8\%) reflecting the vulnerability of hospitalized patients who often have compromised immune systems, underlying conditions, or invasive devices (catheters, ventilators) that increase HAI severity.
    
    \item \textbf{Contamination Decay} (15 seconds): Slow environmental decay reflecting the persistence of healthcare-associated pathogens on surfaces, particularly antibiotic-resistant organisms.
    
    \item \textbf{Immunity Duration} (45 seconds): Extended immunity period modeling post-infection resistance, though notably shorter than lifelong immunity to allow demonstration of re-infection cycles within simulation timeframes.
\end{itemize}

The spatial layout promotes diverse contact patterns: nurses move between wards via the central corridor, creating bridging contacts between patient populations; patients remain largely stationary at beds, representing immobile or mobility-limited individuals; and environmental contamination accumulates around high-traffic areas (nurse stations, corridor intersections), creating persistent transmission reservoirs. This configuration enables study of superspreader dynamics among mobile healthcare workers and the effectiveness of cohorting strategies (isolating infected patients to specific wards).

\subsubsection{School Classroom Scenario}
\label{sec:school-scenario}

The school scenario (\texttt{school\_classroom.json}) models childhood disease transmission (e.g., influenza, chickenpox, measles) in a standard classroom environment with 30 student desks arranged in 5 rows of 6 columns, plus a teacher's area at the front. This scenario represents high-density, high-contact environments where children interact closely during educational activities.

The simulation parameters are tuned to reflect the distinct epidemiological characteristics of school outbreaks:

\begin{itemize}
    \item \textbf{Infection Radius} (120px): Larger transmission radius than hospital settings, reflecting the expansive social interaction patterns of children who engage in close physical contact during play, group activities, and shared workspace use.
    
    \item \textbf{Infection Probability} (0.45 per second): High transmission rate (2.25$\times$ hospital rate) modeling the efficiency of childhood disease spread through direct contact, droplet transmission during close conversation, and shared materials (books, supplies, toys).
    
    \item \textbf{Floor Infection Probability} (0.10): Lower environmental transmission than hospitals, reflecting that while classroom surfaces can harbor pathogens, they typically lack the high bioburden and persistent contamination characteristic of healthcare facilities.
    
    \item \textbf{Recovery Time} (12 seconds): Shorter infection duration than hospital settings, reflecting the typical course of common childhood illnesses (3-5 days for influenza, 5-7 days for chickenpox scaled to simulation time) in otherwise healthy pediatric populations.
    
    \item \textbf{Fatality Probability} (0.02): Very low mortality rate (2\%) reflecting the generally mild course of common childhood diseases in healthy school-age children, though not zero to acknowledge rare complications.
    
    \item \textbf{Contamination Decay} (8 seconds): Faster environmental decay than hospitals, reflecting lower pathogen loads and regular cleaning schedules in school environments.
    
    \item \textbf{Immunity Duration} (20 seconds): Shorter immunity period enabling demonstration of outbreak waves and re-infection dynamics within reasonable simulation timeframes, though many childhood diseases confer longer-lasting immunity in reality.
\end{itemize}

The classroom layout creates distinct transmission dynamics: students seated in adjacent desks experience frequent close-range contact; the teacher acts as a potential superspreader node moving throughout the classroom; and the uniform desk arrangement creates predictable transmission chains along rows and columns. This configuration enables study of desk-spacing interventions, cohorting strategies (keeping groups separate), and the role of high-connectivity individuals (teachers) in outbreak propagation.

\subsubsection{Parameter Tuning Rationale}
\label{sec:parameter-rationale}

The contrast between hospital and school scenarios demonstrates how config-driven parameterization enables modeling of fundamentally different epidemiological contexts without code modification:

\textbf{Transmission Intensity}: School settings exhibit higher infection probability (0.45 vs. 0.20) and larger infection radius (120px vs. 100px) due to the high-contact nature of childhood interactions. Children engage in physical play, share materials, and maintain closer conversational distances than the more structured patient-caregiver interactions in hospitals.

\textbf{Environmental Persistence}: Hospital floors maintain contamination longer (15s decay vs. 8s) and have higher transmission probability (0.20 vs. 0.10), reflecting the presence of antibiotic-resistant organisms and invasive medical procedures that introduce high pathogen loads. Schools, while not sterile, have lower bioburden and benefit from routine cleaning.

\textbf{Disease Severity}: Hospital scenarios have 4$\times$ higher fatality rate (8\% vs. 2\%) reflecting the vulnerability of hospitalized populations with underlying conditions. School outbreaks, while disruptive, typically cause mild illness in healthy children.

\textbf{Temporal Dynamics}: School infections resolve faster (12s vs. 20s) due to robust pediatric immune responses, while hospital infections persist longer due to compromised host defenses and the nature of opportunistic pathogens affecting already-ill patients.

These parameter choices align with epidemiological literature on healthcare-associated infections and childhood disease outbreaks \cite{silva2020covid,kerr2021covasim}, though scaled to demonstration-appropriate timeframes. The config-driven approach enables researchers to adjust these parameters to model specific pathogens (e.g., increasing fatality for COVID-19 in elderly care facilities, increasing transmission for measles in undervaccinated communities) or test intervention scenarios (e.g., reducing infection radius to model mask mandates, increasing contamination decay to model enhanced cleaning protocols).
